\scanpage{9}

\deleted{\uline{Définition} Un \deleted{\unsure{groupe est dit|groupe}} schéma en groupes $X$ \inserted{de type fini} sur $S$ est dit \unsure{tto-ri} de}

\uline{Définition} D'une condition sur un préschéma en \unsure{groupe|groupes} $G$ de \deleted{\inserted{\unsure{prés|plat}} type \unsure{finie|fini}} sur $S$ : \deleted{\unsure{localement|localement noeth}} pour tt schéma local $T$ sur $S$, \deleted{\unsure{il|et}} tout sous-groupe ouvert \deleted{et fermé} $U$ de $G_T$ qui majore la fibre spéciale, on a $U=G_T$.
% slanted
\margin{\uline{Condition R}}
Lorsque \deleted{$G$ \unsure{sur $S$ est|est}} $G$ est de type fini sur $S$ loc. noeth, il suffit de tester aux $T$ spectr. d'un anneau de valuation discrète [complet, corps rés. alg. clos].

On trouve donc

\uline{Corollaire} Soit $\deleted{G}X$ de \deleted{type fini sur présentation}\inserted{type} finie sur $S$ \inserted{loc noeth}, $\deleted{H}Y$ un sous-groupe fermé de, \deleted{présentation finie}. Supposons que $X$ satisfasse R. Alors $\displaystyle\prod_{X/S}Y/X$ est représentable par un sous-préschéma \uline{fermé} de $S$.

\unsure{Exemples|\deleted{Exemples}}

\deleted{Soit}

\deleted{\uline{Corollaire} Soit $G$ schéma en groupes noeth sur $S$ loc. noeth, $H$ un sous-groupe fermé de type fini sur $S$, satisfaisant la condition R. Alors $\underline{\mathrm{Cent}}._{\unsure{G|B}}(H)$ \inserted{(si $G$ séparé sur $S$)} et $\underline{\mathrm{Norm}}_G(H)$ sont représentables par un ss-préschéma fermé de $G$. De $=$ $\mathrm{Transp}_G(H,K)$}
% archive sheet number: 8
